\chapter{Introduction}
\label{ch:introduction}

%Write around four paragraphs establishing the context and motivating your project.

% \section{Related work}

% Discuss related work.

\section{Story}

% One sentence summary of your project. Followed by a short list of concrete objectives:
The restricted three-body system is a model for studying
the motion of an earth-moon satellite or of an asteroid close
enough to the earth to be strongly influenced by the earth
and the sun.

In this system, the two heavy bodies are regarded as revolving
 in fixed orbits about their common centre of mass
while the small body is attracted by the heavy bodies but
never affecting their motion. 
To further simplify the problem, we approximate orbits of the 
heavy bodies as circles so
that we can work in a frame of reference that rotates with
the two heavy bodies. In this frame, the two heavy bodies
are considered as fixed in space with their rotation 
translated into a modification of the equations of gravitational
motion. For simplicity, we scale the units to reduce a number 
of constants to one. We also scale the masses of the two
heavy bodies to $1 - \mu$ and µ and their positions relative to
the moving reference frame by the vectors $\mu e_1$ and $(\mu - 1)e_1$
respectively so that their center of mass is at the origin of
coordinates.

Denote by $u_1$, $u_2$ and $u_3$ the scalar variables representing
the position coordinates of the small body and by $u_4$, $u_5$ and
$u_6$ the corresponding velocity coordinates. The equations of
motion for u(t) are

\begin{align*}
    \begin{cases}
        u_1' = u_4\\
        u_2' = u_5\\
        u_3' = u_6\\
        u_4' = 2u_5+u_1-\frac{\mu(u_1+\mu-1)}{(u_2^2+u_3^2+(u_1+\mu-1)^2)^{3/2}}
        -\frac{(1-\mu)(u_1+\mu)}{(u_2^2+u_3^2+(u_1+\mu)^2)^{3/2}}\\
        u_5' = -2u_4 + u_2 - \frac{\mu u_2}{(u_2^2+u_3^2+(u_1+\mu)^2)^{3/2}}
        -\frac{(1-\mu)u_2}{(u_2^2+u_3^2+(u_1+\mu)^2)^{3/2}}\\
        u_6' = -\frac{\mu u_3}{(u_2^2+u_3^2+(u_1+\mu-1)^2)^{3/2}}-\frac{(1-\mu)u_3}{(u_2^2+u_3^2+(u_1+\mu)^2)^{3/2}}
    \end{cases}
\end{align*}

Planar motion is possible; that is, solutions which satisfy $u_3 = u_6 = 0$ at all times. One of these is shown in the
LHS plot of Figure 10.14, with the values of $(u_1, u_2)$ plotted as the orbit evolves. The heavier mass is at the point
$(−\mu, 0)$ and the lighter mass is at $(1 - \mu, 0)$, where $(0, 0)$ is
marked 0 and (1, 0) is marked 1. For this calculation the
value of $\mu = 0.012277471$ was selected, corresponding to the
earth-moon system. The initial values for this computation
were
\begin{align*}
    u(0) = (0.994, 0, 0, 0, -2.0015851063790825224, 0) && (10.190)
\end{align*}
and the period was $T_1 =17.06521656015796$.

The RHS plot of Figure 10.14 illustrates another solution
of (10.189) with a different initial condition
\begin{align*}
    u(0) = (0.879779227778, 0, 0, 0, −0.379677780949, 0) && (10.191)
\end{align*}
and a different period $T_2 =19.140540691377$.

\includegraphics*[scale = 0.8]{../figures/f10.14.png}

\section{编译}

注：由于在确定满足误差小于一定值的最小cpu时间时，采用了二分搜索法，以及控制自适应步长算法的总步数时采用了遍历搜索方法，
导致运行完整的程序需要的时间较长，建议注释一部分，运行一部分。在自行测试代码时，如果不使用求解器工厂获得求解器，请注意使用
智能指针以便及时释放内存。

\subsection{使用xmake编译}
\begin{verbatim}
$ xmake build ivp //生成静态库
$ xmake build test //编译测试程序
$ xmake run test //运行测试程序
$ xmake run plot //将测试程序的输出绘制成report需要的图片
$ xmake build report //编译latex文件程序报告
$ xmake clean //清除所有生成的文件
$ xmake build all //按顺序运行以上除了clean的所有
\end{verbatim}


\subsection{使用make编译}